\subsection{完全交理想与正则环}

命题 2.--- 设 A 为 Noether 环。下列条件等价：

\begin{enumerate}
\def\labelenumi{(\roman{enumi})}
\item
  A 是正则的；
\item
  A 的每个极大理想都完全割；
\item
  A 的每个使 A/J 正则的理想 J 都完全割。
\end{enumerate}

(i)⇒(iii)：设环 A 正则；设 J 为 A 的理想使 Λ/J 正则，并设 p 为 A 的一个包含 J 的素理想。则局部环 \(A_{p}\) 与 \(A_{p}/J_{p}\) 都正则，于是 \(J_{p}\) 由一个对 \(A_{p}\) 完全割的序列生成 (VIII, § 5, 第 3 节, 推论 2 与命题 2)，这就意味着 J 在 p 处完全割。

\begin{enumerate}
\def\labelenumi{(\roman{enumi})}
\setcounter{enumi}{2}
\tightlist
\item
  (ii)：这是显然的，因为域是正则环。
\end{enumerate}

(ii)⇒(i)：设 m 为 A 的一个极大理想；在假设 (ii) 下，极大理想 mA \(_{m}\)（Λ \(_{m}\) 中的）由一个对 A \(_{m}\) 完全割的序列生成，因此 Λ \(_{m}\) 正则 (VIII, § 5, 第 2 节, 定理 1)。于是 A 正则。

命题 3.--- 设 A 为 Noether 环，J 为 A 的理想且 A/J 正则。下列条件等价：

\begin{enumerate}
\def\labelenumi{(\roman{enumi})}
\item
  理想 J 完全割；
\item
  对 A 的每个包含 J 的素理想（相应地极大理想）p，环 A \(_{p}\) 都正则；
\item
  A 关于 J-进拓扑的分离完备化是正则的。
\end{enumerate}

由定理 1，条件 (i) 意味着：对 A 的每个包含 J 的素理想（相应地极大理想）p，理想 \(J_{p}\)（局部环 \(A_{p}\) 中的）由一个对 \(A_{p}\) 完全割的序列生成。由于按假设 \(A_{p}/J_{p}\) 正则，后一条件等价于 \(A_{p}\) 正则 (VIII, § 5, 第 3 节, 命题 2 及其推论 2)；这就证明了 (i) 与 (ii) 的等价性。(ii) 与 (iii) 的等价性由 § 4, 第 2 节, 命题 4 的推论 3 得出。

命题 4.--- 设 A 为正则环，J 为 A 的理想，\(A_{0}\) 为 A 的子环且典范同态 \(A_{0} \rightarrow A/J\) 是双射的。

\begin{enumerate}
\def\labelenumi{\alph{enumi})}
\item
  理想 J 完全割，\(A_{0}\) -模 \(J/J^{2}\) 是有限型投射的，且环 \(A_{0}\) 正则。
\item
  设 \(\varphi: J/J^{2} \to J\) 为一个 \(A_{0}\) -线性截面（对典范满射 \(J \to J/J^{2}\)）。\(A_{0}\) -代数同态 \(\mathsf{S}_{A_{0}}(J/J^{2}) \longrightarrow A\)（它扩张 \(\varphi\)）扩张为分次环 \(\mathsf{S}_{A_{0}}(J/J^{2})\) 的完备化到 \(A\) 关于 \(J\) -进拓扑的分离完备化上的一个同构。
\end{enumerate}

\begin{enumerate}
\def\labelenumi{\alph{enumi})}
\item
  设 \(\mathfrak{p}\) 为 A 的一个包含 J 的素理想。我们有 \(\mathfrak{p} = (\mathfrak{p} \cap A_0) \oplus J\)，从而 \(\mathfrak{p}^2 = (\mathfrak{p} \cap A_0)^2 \oplus \mathfrak{p}J\)，故 \(\mathfrak{p}^2 \cap J = \mathfrak{p}J\)。设 \(i\) 为 \(J\) 到 \(\mathfrak{p}\) 中的典范单射。由前所述，映射 \(i \otimes 1_{A/p}: J \otimes_A A/p \longrightarrow \mathfrak{p} \otimes_A A/p\) 是单射的，对 \(i_{\mathfrak{p}} \otimes 1_{\kappa(\mathfrak{p})}: J_{\mathfrak{p}} \otimes_{A_{\mathfrak{p}}} \kappa(\mathfrak{p}) \longrightarrow \mathfrak{p}A_{\mathfrak{p}} \otimes_{A_{\mathfrak{p}}} \kappa(\mathfrak{p})\) 亦然。Nakayama 引理蕴含：理想 \(J_{\mathfrak{p}}\)（\(A_{\mathfrak{p}}\) 中的）由正则局部环 \(A_{\mathfrak{p}}\) 的一个坐标系的一个子集生成。根据 VIII, § 5, 第 3 节的命题 2，环 \(A_{\mathfrak{p}}/J_{\mathfrak{p}}\) 正则且理想 \(J_{\mathfrak{p}}\) 完全割。于是 \(J\) 完全割，环 \(A_0\) 同构于 \(A/J\) 故正则，且由定理 1（第 2 节），\(A_0\) -模 \(J/J^2\) 是有限型投射的。
\item
  由定理 1，典范同态 \(\beta_{\mathrm{J}}: \mathsf{S}_{\mathrm{A}_{0}}(\mathrm{J}/\mathrm{J}^{2}) \longrightarrow \mathrm{gr}_{\mathrm{J}}(\mathrm{A})\) 是双射的。设 \(f: \mathsf{S}_{\mathrm{A}_{0}}(\mathrm{J}/\mathrm{J}^{2}) \longrightarrow \mathrm{A}\) 为 \(\mathrm{A}_{0}\) -代数同态，它扩张 \(\mathrm{A}_{0}\) -线性映射 \(\varphi: \mathrm{J}/\mathrm{J}^{2} \longrightarrow \mathrm{J}\)。若给 A 赋以 J-进滤链，并给 \(\mathsf{S}_{\mathrm{A}_{0}}(\mathrm{J}/\mathrm{J}^{2})\) 赋以与其分次相伴的滤链，则 \(\beta_{\mathrm{J}}\) 等同于由 \(f\) 通过取相伴分次对象得到的同态。于是断言 b) 由 III, § 2, 第 8 节, 定理 1 的推论 3 得出。
\end{enumerate}

